% !TEX root = ../fastmat.tex

% Introduce the scientific background and the motivation for developing the software.
%
% Explain why the software is important, and describe the exact (scientific) problem(s) it solves.
%
% Indicate in what way the software has contributed (or how it will contribute in the future) to the process of scientific discovery; if available, this is to be supported by citing a research paper using the software.
%
% Provide a description of the experimental setting (how does the user use the software?).
%
% Introduce related work in literature (cite or list algorithms used, other software etc.).

\begin{itemize}
\item linear mappings are present in applied math, physics, engineering, signal processing and data science
\item represented as matrices when acting between finite dimensional spaces
\item many mappings are not arbitrary but follow a certain regularity, logic or physics: convolution, differentiation, integration
\item they are structured: fewer degrees of freedom
\item naturally: this structure is also present in our minds, our notation, our thinking
\item so the represenation as a matrix becomes wasteful and cumbersome
\item also: applying the mapping is not a matrix-vector product anymore, but an algorithm on its own, specifically tailored to the mapping
\item we have to drop the represenation as a matrix
\item we want to combine the rectangle of numbers and the specifically tailored algorithm into one concept
\item many algorithms only need the mapping and not the whole matrix (cg, bicgstab, lancosz iteration, ista, omp)
\end{itemize}
